%% ====================================================================
\section{Implementation Architecture}
\label{sec:architecture}
%% ====================================================================

\subsection{File Structure}
\label{sec:files}

The project is organized into four directories: \texttt{src/} (core computation),
\texttt{test/} (validation), \texttt{gpu/} (CUDA kernels), and \texttt{data/}
(reference tables).

\paragraph{Core source files (\texttt{src/}).}

\begin{longtable}{@{}p{4.2cm}p{8.5cm}@{}}
\toprule
\textbf{File} & \textbf{Purpose} \\
\midrule
\endhead
\texttt{dirichlet.h/.cpp}      & Partial Dirichlet series $S_N(s)$ and $\eta_N(s)$ \\
\texttt{partition.h/.cpp}      & Partition function $Z(\beta) = \sum n^{-\beta}$ with precomputation \\
\texttt{theta.h/.cpp}          & Riemann-Siegel theta function $\vartheta(t)$; Stirling expansion with error control \\
\texttt{system1.h/.cpp}        & $\mathcal{L}(\beta,t)$: accumulated phase factor (System~1) with alternating series \\
\texttt{system2.h/.cpp}        & $\mathcal{G}(\beta,t)$: Loschmidt amplitude (System~2) with $Z$-function connection \\
\texttt{rate\_function.h/.cpp} & Rate function $\mathcal{F}(\beta,t)$ with underflow guard and sentinel detection \\
\texttt{zeros.h/.cpp}          & Zero detection: bisection, Gram points, Turing zero-count check \\
\texttt{euler\_maclaurin.h/.cpp} & Euler-Maclaurin tail correction; Bernoulli number computation \\
\texttt{kahan.h}               & Kahan compensated summation \\
\texttt{main.cpp}              & Entry point; CLI argument parsing \\
\texttt{config.h}              & Compile-time and runtime configuration \\
\bottomrule
\end{longtable}

\paragraph{Test suite (\texttt{test/}).}

\begin{longtable}{@{}p{5cm}p{7.7cm}@{}}
\toprule
\textbf{File} & \textbf{Purpose} \\
\midrule
\endhead
\texttt{test\_dirichlet.cpp}        & Unit tests for $S_N$ and $\eta_N$ against \texttt{mpmath} \\
\texttt{test\_theta.cpp}            & Unit tests for $\vartheta(t)$ \\
\texttt{test\_known\_zeros.cpp}     & Verification against first 100 known zeros \\
\texttt{test\_convergence.cpp}      & Convergence rate $O(N^{-\beta})$ verification \\
\texttt{test\_off\_line.cpp}        & Off-critical-line scan \\
\texttt{test\_rate\_function.cpp}   & Rate function value checks (Controls 1--2) \\
\texttt{test\_cross\_validation.cpp}& System~1 vs System~2 consistency \\
\texttt{reference\_values.h}        & High-precision reference values from \texttt{mpmath} \\
\bottomrule
\end{longtable}

\paragraph{GPU kernels (\texttt{gpu/}).}

\begin{longtable}{@{}p{4.2cm}p{8.5cm}@{}}
\toprule
\textbf{File} & \textbf{Purpose} \\
\midrule
\endhead
\texttt{sweep.cu}           & CUDA kernel for $(\beta,t)$ grid evaluation \\
\texttt{gpu\_dirichlet.cuh} & Device-side Dirichlet series \\
\texttt{gpu\_config.cuh}    & GPU-specific configuration \\
\bottomrule
\end{longtable}

\paragraph{Data and output.}
\texttt{data/zeros\_known.txt} contains the first $100+$ known nontrivial zeta zeros
(sourced from the LMFDB~\cite{lmfdb} and Odlyzko~\cite{odlyzko1989}).
The \texttt{results/} directory is populated at runtime.


\subsection{Language Choices}
\label{sec:languages}

\paragraph{Core computation: C++17.}
Chosen for:
\begin{itemize}
  \item Direct control over floating-point arithmetic (IEEE~754 compliance via compiler
        flags \texttt{-ffp-contract=off}, \texttt{-fno-fast-math}).
  \item Efficient memory management for large arrays.
  \item Compatibility with CUDA (which requires C++ host code).
  \item No garbage collection pauses during long computations.
\end{itemize}

\paragraph{GPU acceleration: CUDA C++.}
Chosen for:
\begin{itemize}
  \item Native support on the target hardware (NVIDIA RTX~2060).
  \item Fine-grained control over thread/block organization.
  \item Access to hardware special function units for \texttt{sin}/\texttt{cos}/\texttt{exp}.
\end{itemize}

\paragraph{Test validation: Python/mpmath (external).}
Reference values are precomputed using \texttt{mpmath} with $50$-digit precision and
hardcoded into the C++ test suite. This removes any runtime dependency on Python.


\subsection{Dependencies}
\label{sec:deps}

The implementation is designed to minimize external dependencies:
\begin{itemize}
  \item \textbf{C++ standard library} (C++17): containers, algorithms, I/O,
        \texttt{<cmath>}.
  \item \textbf{CUDA Toolkit} ($\ge 11.0$): for GPU kernels. Optional; the code compiles
        and runs in CPU-only mode without CUDA.
  \item \textbf{No other external dependencies.} No Boost, no Eigen, no BLAS/LAPACK, no
        numerical libraries. All required mathematical functions (Stirling expansion,
        Kahan summation, Bernoulli numbers, bisection) are implemented from scratch for
        full transparency and auditability.
\end{itemize}


\subsection{Build System}
\label{sec:build}

A single Makefile with the following targets:

\begin{center}
\begin{tabular}{@{}ll@{}}
\toprule
\textbf{Command} & \textbf{Action} \\
\midrule
\texttt{make}       & Build CPU-only version \\
\texttt{make gpu}   & Build with CUDA GPU support \\
\texttt{make test}  & Build and run all unit tests \\
\texttt{make clean} & Remove build artifacts \\
\bottomrule
\end{tabular}
\end{center}

Compiler flags:
\begin{itemize}
  \item \texttt{-std=c++17 -O2 -Wall -Wextra -Wpedantic}
  \item \texttt{-ffp-contract=off} (prevent fused multiply-add, which can change results)
  \item \texttt{-fno-fast-math} (ensure IEEE~754 compliance)
  \item For CUDA: \texttt{-arch=sm\_75} (Turing architecture, RTX~2060)
\end{itemize}
